\scanpage{10}
\mspage{2}
qui est un 2-\inserted{pseudo-}torseur sous une 2-Pic-catégorie, et ce 2-torseur est 2-équivalent à celui des hom de 2-gerbes sur $B_G$ liées par $M'_{B_G}$ de
\margin{$M_{B_G}\unsure{\relbar|\longrightarrow}M'_{B_G}$\\[3pt]$B_G\xrightarrow{B(u)}B_G$}
$\widetilde{\unsure{\mathscr P|P}}\otimes_{M_{B_G}}M'_{B_G}$ dans $B_{\unsure{u}}^*(\widetilde{\mathscr P}')$. Les invariants de la 2-Pic-catégorie pertinente sont $H^2(G,M'),H^1(G,M'),H^0(\unsure{C|G},M')$. L'obstruction à trouver un objet dans le 2-ps.-torseur est dans $H^3(G,M')$, et cela se calcule comme différence de deux classes, images canoniques des éléments qu'on devine dans $H^3(C,M)$ et $H^3(C',M')$.

4. Généralisation aux Gr-champs sur un topos $S$. \deleted{\unsure{On les|Bien sûr}} Si $G$ (groupe sur $S$) et $M$ ($G$-Module sur $S$) sont donnés, on interprète tout en termes de 2-gerbes ``relatives'' sur $B_G$ mod $S$ liées par $M_{B_G}$, i.e. 2-gerbes sur $B_G$ trivialisées sur $S$ (par $S\simeq B_e\longrightarrow B_G$)

Il faudrait le mettre ``en \unsure{série|séries|scène}'' avec l'histoire \unsure{des|de} \uline{1-gerbes} \unsure{relative|relatives} sur $B_C$, correspondant aux \uline{extensions} de $G$ par $M$. Il faudrait traiter en $=$ temps le cas d'un $M$ non commutatif\ldots
% archive sheet number: 9
